\section{Relation to prior work and the scope of the contributions}
\label{sec:prior-work}

The mechanisms behind the proofs have substantial histories. The purpose
of this comparison is to identify the quantitative conclusions and
integrations proved here, while attributing the tools on which they rest.
It is a statement-level comparison, not a claim that a keyword search
establishes priority.

\paragraph{Elimination and symbolic ordinary inputs.}
Bucket elimination and its combinations with conditioning are established
methods \citep{Dechter1999}. The generalized distributive law applies
marginalization to products in a commutative semiring \citep{Aji2000}.
Taking that semiring to be Boolean functions of $x$, with pointwise OR
and AND, explains why ordinary inputs can remain symbolic. Storing each
function as a circuit pointer makes a semiring operation cost one gate.
Thus symbolic elimination is an application of an existing principle.
The work in Section~\ref{sec:consistency-graph} is its concrete
$B_2$ implementation, consistency budget, cubic reduction, and exact
exponential-rate accounting.

\paragraph{Knowledge compilation has stronger target restrictions.}
DNNF supports efficient existential forgetting
\citep{DarwicheMarquis2002}. Bounded-treewidth Boolean circuits can be
compiled into structured deterministic DNNFs with singly exponential
dependence on width \citep{Amarilli2020}. Here ``deterministic'' means
that alternatives at an OR gate are mutually exclusive. An ordinary
deterministic Boolean circuit has no such restriction. Forgetting may
lose that property while still yielding a perfectly valid target for
this note. Results retaining stronger structure through quantification,
such as those of \citet{CapelliMengel2019}, therefore solve a different
representation problem. Their costs and lower bounds are not lower
bounds on unrestricted $\C(g)$.

Full circuit treewidth, treewidth after deleting input vertices, and
the consistency graph used here are also different parameters. Prior
bounded-treewidth circuit-SAT work explicitly uses an input-deleted
graph \citep{Lokshtanov2018}. Morizumi's structural simulation uses
layered circuit width and gives size
$2^{O((w+\log s)\log s)}$ \citep[Theorem 4]{Morizumi2018}; that width
is not interchangeable with our $\kappa$ or cutwidth.

\paragraph{Graph width and tensors.}
The pathwidth theorem supplying the $1/6$ coefficient is due to
\citet{FominHoie2006}. The conversion between pathwidth and subcubic
cutwidth is also known through graph-search relations; an explicit
modern restatement appears in the proof of Lemma 16 of
\citet{Bodlaender2025}. Our direct median proof makes the additive
constant and resulting exponential base transparent.
Width-controlled tensor contraction is established
\citep{MarkovShi2008}, including Boolean checking circuits and
counting applications \citep{Johnson2013}. The specific quantitative
chain here is
\[
 h+b+\kappa\le q+1,
 \qquad N\le2\kappa,
 \qquad w\le N/6+o(N),
 \qquad \C(g)\le p+O(q2^w).
\]
Together with enumeration and nonuniform synthesis, it yields the
$1/5$ projection rate. None of the inspected statements asserts this
exact combined formula. This is the boundary of the comparison, not
a proof that no equivalent earlier result exists.

There are also direct antecedents for the elementary tensor steps.
\citet{BiamonteMortonTurner2015} expand each COPY tensor into two
consistent choices and evaluate the resulting trees, obtaining a
bound polynomial in the representation size times $2^c$ for $c$
COPY tensors in their canonical representation. This is the mechanism
behind the forest-cut baseline. Its mixed-input specialization here
leaves ordinary inputs symbolic and gives $r\le s+1-a-b$.
\citet{BeaudrapKissingerMeichanetzidis2021} explicitly switch between
decision and counting semirings and study exact tensor rewrites.
\citet{Kourtis2019} give planar contraction guarantees and empirical
studies of general cubic instances. The latter experiments do not
supply a worst-case gate-count exponent.
\citet{DudekDuenasOsorioVardi2020} give exact weighted-counting tensor
representations and tree factorizations of suitable high-degree tensors.
Maximum intermediate rank controls storage; our sequential frontier
calculation separately charges the operations needed for the time bound.
These sources reinforce that
the width-and-cycle accounting, rather than the general use of tensors
or COPY expansion, is the quantitative step to assess here.

Appendix~\ref{sec:aggregation} likewise uses established polynomial
identity testing and algebraic matching. Its transfer statement makes
precise how a cheaper weighted evaluator would imply a smaller exact
projection circuit, while unary tags preserve the cycle budget.
It is an avenue for improving the main bound, not an additional
improvement of its exponent.

\paragraph{Coverage, penalties, and sparse correction.}
The function $H\mapsto\sum_vN_v(H)+N_{\mathrm{out}}(H)$ is a
monotone submodular cost: each term counts covered restriction keys.
Selecting witness columns with penalties for uncovered positive rows
therefore fits an existing submodular-cost set-cover framework
\citep{Wang2015}. No approximation guarantee from that framework
is imported here, and an algorithm polynomial in an explicit incidence
matrix need not be polynomial in a succinct verifier circuit.

Optimizing only minimum row mass has an elementary fractional-cover
interpretation. If $g$ is nonzero and $\tau^*$ is its fractional
witness-cover number, then
\[
 \max_\mu\min_{x:g(x)=1}\mu(W_x)=1/\tau^*.
\]
Indeed, divide a distribution by its minimum row mass to obtain a
fractional cover, and normalize a fractional cover by its total weight
for the converse. The occupancy objective additionally prices
collisions on internal supports. The examples in
Section~\ref{sec:row-mass-gap} prove that even identical row-mass
vectors leave an exponential ambiguity in that price.

Circuit synthesis for functions with few positive inputs is likewise
classical; see \citet{Redkin2004,Redkin2021}. The block construction
in Lemma~\ref{lem:sparse} is a convenient all-range majorant, not a
uniformly optimal sparse-synthesis theorem. In particular, the cited
literature gives linear-size bounds in some ranges with a logarithmic
number of positive inputs. The contribution of
Theorem~\ref{thm:occupancy} is the joint inequality using the same
sample for shared local computation and exact sparse correction,
together with the distribution tradeoffs and hardness consequences
proved from it.

\paragraph{The uniform counting consequence and its comparison.}
Corollary~\ref{cor:gate-counting} makes the all-quantified consequence
explicit. It proves count-preserving reductions, polynomial bit costs,
and a fixed-$\varepsilon$ running-time bound with exponential rate
approaching $1/4$ in gate count. This uniform consequence is separate
from the synthesis step in the nonuniform projection theorem.
Its proof does not establish a claim to a previously best SAT algorithm.

The inspected branching literature records Nurk's general
Circuit-SAT bound $O(2^{0.4058s})$ \citep{Nurk2009} and Savinov's
improvement to $O(2^{0.389667s})$, presented by
\citet{Lialina2020}, and the nontrivial $B_2$ counting threshold
below $3n$ gates of \citet{GKST2018}. Width-based dynamic programming
uses exponential space, whereas depth-first branching can have a
different space cost; width-based SAT time--space tradeoffs are
studied explicitly by \citet{Allender2014}. The differing space
costs do not, by themselves, establish whether the numerical
consequence is new or explain its absence from particular benchmark
discussions. For sufficiently small fixed $\varepsilon$, the proved
counting exponent is smaller than these particular displayed decision
exponents. This arithmetic comparison does not establish historical
priority or the best unrestricted-space running time. The precise
combined result may be a previously known consequence of width-based
methods; the source comparison has not resolved that question.

\paragraph{What is established here.}
The displayed proofs establish the consistency-graph simulation and
its $1/5$ rate, the necessary critical structure, the exact kernel
realization with nonminimal verifiers, the occupancy--alteration bound,
the quantitative sparse-fiber population, the all-quantified counting
consequence, and the worked separations between compilation
guarantees. The realization begins with the classical $st$-numbering
property and separately proves its repeated-input repair budget and
exact retained graph. It does not make a claim about the kernels of
minimum verifiers. The ingredients attributed above are established prior
methods. Priority for the precise combined statements is unestablished
beyond the documented comparison. No unrestricted exponential lower
bound, exact optimal exponent, or resolution of the remaining hard
range is claimed.
